\section{Fix counts, capacities and ownership}
The four whole-buffer entries receive \texttt{a0}: source pointer, \texttt{a1}: unsigned
64-bit input count, \texttt{a2}: destination pointer and \texttt{a3}: unsigned 64-bit
output capacity. The scalar and MSA versions have the same results and failure rules.
All pointers may be byte-aligned. Source and destination are disjoint for encoding;
decoding permits either disjoint buffers or exactly the same starting address. Other
partial overlap is outside this project.

\tagpdfsetup{table/header-rows={1}}
\begin{longtable}{p{.30\textwidth}p{.34\textwidth}p{.24\textwidth}}
Entry & Selected input & Selected output \\
\hline
\texttt{b64\_encode\_scalar} & At most 4096 binary bytes & \(4\lceil n/3\rceil\) ASCII bytes \\
\texttt{b64\_encode\_msa} & The same binary domain & The same encoded bytes \\
\texttt{b64\_decode\_scalar} & At most 5464 canonical text bytes & \(3(n/4)-p\) binary bytes \\
\texttt{b64\_decode\_msa} & The same canonical text domain & The same decoded bytes \\
\texttt{b64\_encode12\_msa} & Exactly twelve binary bytes & Exactly sixteen ASCII bytes \\
\texttt{b64\_decode16\_msa} & Sixteen alphabet bytes, no padding & Exactly twelve binary bytes \\
\end{longtable}

Here \(p\) is zero, one or two final padding characters, after the text has been
validated. The kernels use \texttt{a0} for source and \texttt{a1} for destination.
They require those exact mapped intervals and return success; the decoder kernel's
input must already be known valid and unpadded. Its buffers may be disjoint or exactly
in place. The encoder kernel's buffers are disjoint. The kernels do not infer capacity
or accept a partial group.

Whole-buffer success returns \texttt{v0}=0 and the exact output byte count in
\texttt{v1}. The kernels return counts sixteen and twelve respectively. No entry adds
a terminating zero. On a whole-buffer error, \texttt{v1} is zero and caller memory is
unchanged. Status one rejects an input count above the entry's bound, before pointer
access. Status two rejects malformed decoder input; a count not divisible by four is
rejected before reading source. Status three reports insufficient output capacity.

An empty whole-buffer request returns success with count zero before accessing either
pointer, even if both are invalid and capacity is zero. For nonempty encoding, the
required size is known from the input count, so check capacity before reading source.
For nonempty decoding, validate text first, then check capacity. Malformed text therefore
returns status two even when capacity is also too small. A capacity of \(2^{63}\) is a
large unsigned value, not a negative buffer size. Only the bytes actually consumed or
produced must be mapped.

Entries preserve every bit of \texttt{s0} through \texttt{s7}, \texttt{fp} and
\texttt{sp}, return through the full supplied \texttt{ra}, and perform no console I/O.
The caller supplies an eight-byte-aligned stack. Temporary scalar and vector registers are volatile in this teaching interface. The
illustrated kernels use \texttt{w0} through \texttt{w7} and also change their
floating-point aliases. Read only selected input or output intervals among caller-owned objects,
and write only the selected output and private storage. Constants and private workspace
are permitted. Exact in-place decoding changes the output prefix of its input object;
bytes beyond that prefix remain unchanged.
